% META: {"id": "1NSI-TABLES-RE-C4", "chapitre": "1NSI-TABLES", "type_objet": "remediation", "capacites": ["P-TABLE-04"], "capacites_codes": ["C4"], "mode_creation": "ex_nihilo", "sources_inspiration": [], "status": "needs_review"}

%---------------------------------------------------------------
% FICHE DE REMEDIATION — C4 : la fusion ne garde que le domaine commun
%---------------------------------------------------------------

\textbf{Ce qui bloque.} Tu attends de la fusion qu'elle garde toutes les lignes des deux
tables, ou qu'elle « complète » celles qui n'ont pas de correspondance. Par défaut, seules les
valeurs de la clé présentes \emph{dans les deux} tables produisent une ligne. Et quand un
nom de colonne existe des deux côtés, la fusion s'arrête sur une erreur au lieu d'écraser
une valeur en silence.

\textbf{À retenir.} Avant de fusionner, écrire les deux domaines de la clé et leur
intersection : c'est elle qui donne le nombre de lignes. Puis vérifier que les colonnes hors
clé portent des noms différents.

\begin{exercice}{1NSI-TAB-RE-C4-EX1}{1}{12}
Le lycée dispose de deux tables :
\begin{python}
eleves = [
    {"nom": "Amine", "classe": "1A"},
    {"nom": "Lina", "classe": "1B"},
    {"nom": "Omar", "classe": "1A"},
    {"nom": "Sara", "classe": "1B"},
]
clubs = [
    {"nom": "Lina", "club": "robotique"},
    {"nom": "Sara", "club": "theatre"},
    {"nom": "Yanis", "club": "echecs"},
]
\end{python}
Un élève prévoit que \lstinline{fusionner(eleves, clubs, "nom")} renverra sept lignes, «
quatre élèves plus trois inscriptions ».
\begin{enumerate}
  \item Écrire le domaine de la clé \lstinline{nom} dans chaque table, puis leur
        intersection. Combien de lignes la fusion renvoie-t-elle réellement ? Les donner.
  \item Amine et Omar ne sont dans aucun club, Yanis n'est pas dans \lstinline{eleves}.
        Pourquoi aucun des trois n'apparaît-il ? Ces cas sont-ils des erreurs ?
  \item On remplace \lstinline{clubs} par une table où chaque ligne porte aussi une colonne
        \lstinline{classe}. Que fait \lstinline{fusionner} ? Pourquoi ce comportement
        vaut-il mieux qu'un résultat silencieux ?
\end{enumerate}
\end{exercice}

% BEGIN-VERIFY
% def fusionner(table1, table2, cle):
%     colonnes1 = {nom for ligne in table1 for nom in ligne if nom != cle}
%     colonnes2 = {nom for ligne in table2 for nom in ligne if nom != cle}
%     if not colonnes1.isdisjoint(colonnes2):
%         raise ValueError("les colonnes hors cle doivent etre disjointes")
%     index2 = {ligne[cle]: ligne for ligne in table2}
%     fusion = []
%     for ligne1 in table1:
%         if ligne1[cle] in index2:
%             complement = {
%                 k: v for k, v in index2[ligne1[cle]].items() if k != cle
%             }
%             fusion.append({**ligne1, **complement})
%     return fusion
% eleves = [
%     {"nom": "Amine", "classe": "1A"},
%     {"nom": "Lina", "classe": "1B"},
%     {"nom": "Omar", "classe": "1A"},
%     {"nom": "Sara", "classe": "1B"},
% ]
% clubs = [
%     {"nom": "Lina", "club": "robotique"},
%     {"nom": "Sara", "club": "theatre"},
%     {"nom": "Yanis", "club": "echecs"},
% ]
% commun = {ligne["nom"] for ligne in eleves} & {ligne["nom"] for ligne in clubs}
% assert commun == {"Lina", "Sara"}
% resultat = fusionner(eleves, clubs, "nom")
% assert resultat == [
%     {"nom": "Lina", "classe": "1B", "club": "robotique"},
%     {"nom": "Sara", "classe": "1B", "club": "theatre"},
% ]
% assert len(resultat) != 7
% try:
%     fusionner(eleves, [{"nom": "Lina", "classe": "1B", "club": "robotique"}], "nom")
% except ValueError as erreur:
%     assert str(erreur) == "les colonnes hors cle doivent etre disjointes"
% else:
%     raise AssertionError("une collision de colonnes doit etre refusee")
% END-VERIFY
